% ======================================================================
% Contribution 5: Axiomatic Derivation of Coupling Form Optimality
% ======================================================================

\section{Optimality of the Coupling Forms}
\label{sec:coupling}

The companion TMLR paper introduces two coupling forms for
sigma-targeting curricula: additive coupling
$\mathcal{L}_{\mathrm{mod}} = \mathcal{L}_{\mathrm{task}} (1 + C(1 - \sigma_A))$
and multiplicative coupling
$\mathcal{L}_{\mathrm{mod}} = \mathcal{L}_{\mathrm{task}} (1 + C \sigma_A)$,
with the $\Psi_A$ intersection activation taking the geometric mean form
$\Psi_A(d_1,d_2) = \Psi_0 \cdot \phi(d_1,d_2) \sqrt{\sigma_A(d_1) \sigma_A(d_2)}$.
However, no proof is given that these forms are optimal or unique. In this
section, we derive the coupling forms from first principles using an
axiomatic framework, proving that the geometric-mean $\Psi_A$ and the
multiplicative coupling are the unique functions satisfying natural
physical and information-theoretic axioms.

\subsection{Axioms for Intersection Activation}

Let $\Psi_A(d_1, d_2)$ be the rate at which the agent discovers useful
intersections between domains $d_1$ and $d_2$. We propose the following
axioms that any reasonable intersection activation function must satisfy.

\begin{axiom}[Symmetry]
\label{ax:symmetry}
$\Psi_A(d_1, d_2) = \Psi_A(d_2, d_1)$.
Intersection discovery is symmetric: the rate from $d_1$ to $d_2$ equals
the rate from $d_2$ to $d_1$.
\end{axiom}

\begin{axiom}[Monotonicity in Schema Coherence]
\label{ax:monotonicity}
$\Psi_A$ is strictly increasing in both $\sigma_A(d_1)$ and $\sigma_A(d_2)$:
\[
\frac{\partial \Psi_A}{\partial \sigma_A(d_1)} > 0,
\qquad
\frac{\partial \Psi_A}{\partial \sigma_A(d_2)} > 0.
\]
Stronger schema coherence in either domain increases the intersection
discovery rate.
\end{axiom}

\begin{axiom}[Zero Coherence $\Rightarrow$ No Transfer]
\label{ax:zero-coherence}
\leavevmode\newline
If $\sigma_A(d_1)=0$ or $\sigma_A(d_2)=0$, then
$\Psi_A(d_1,d_2)=0$. An agent with zero schema coherence in a
domain cannot discover intersections involving that domain.
\end{axiom}

\begin{axiom}[Invariance Under Independent Scaling]
\label{ax:scale-invariance}
For any $\lambda > 0$, scaling both schema coherences by
$\lambda$ should preserve the relative activation rate:
\[
\frac{\Psi_A(d_1, d_2; \lambda \sigma_1, \lambda \sigma_2)}
{\Psi_A(d_1, d_2; \sigma_1, \sigma_2)}
= g(\lambda)
\]
for some function $g$ independent of $\sigma_1, \sigma_2$. This
captures the idea that the functional form is homogeneous.
\end{axiom}

\begin{axiom}[Normalization]
\label{ax:normalization}
When $\sigma_A(d_1) = \sigma_A(d_2) = 1$ (maximal coherence in both
domains) and $\phi(d_1, d_2) = 1$ (identical domains), the activation
rate is $\Psi_A = \Psi_0$ (the base rate).
\end{axiom}

\begin{theorem}[Uniqueness of the Geometric-Mean Form]
\label{thm:psi-unique}
Under Axioms~\ref{ax:symmetry}--\ref{ax:normalization}, the intersection
activation function must take the form
\[
\Psi_A(d_1, d_2) = \Psi_0 \cdot \phi(d_1, d_2) \cdot
\sqrt{\sigma_A(d_1) \sigma_A(d_2)}.
\]
This form is unique up to the scalar multiplier $\Psi_0$.
\end{theorem}

\begin{proof}
By Axiom~\ref{ax:scale-invariance}, the ratio
$\Psi_A(\lambda \sigma_1, \lambda \sigma_2) / \Psi_A(\sigma_1, \sigma_2)$
depends only on $\lambda$. Applying the functional equation twice gives
$g(\lambda \mu) = g(\lambda) g(\mu)$ for $g(\lambda) \defeq \Psi_A(\lambda, \lambda) / \Psi_A(1, 1)$.
Since $\Psi_A$ is continuous (it is a $C^1$ function of its arguments
given the smoothness of the ODE system), $g$ is continuous, and the only
continuous solutions of the Cauchy exponential equation are
$g(\lambda) = \lambda^\alpha$ for some $\alpha \in \R$. Hence $\Psi_A$
is homogeneous of degree $\alpha$:
\[
\Psi_A(\lambda \sigma_1, \lambda \sigma_2) = \lambda^\alpha \Psi_A(\sigma_1, \sigma_2).
\]

By Axiom~\ref{ax:zero-coherence}, $\Psi_A$ vanishes when either argument
is zero. This forces $\alpha > 0$, because $\Psi_A(\lambda \sigma_1, 0) = 0$
for all $\lambda$ and the only homogeneous function that vanishes on the
coordinate axes for all scales has positive degree.

By Axiom~\ref{ax:symmetry}, $\Psi_A$ is symmetric. The general symmetric
homogeneous function of degree $\alpha$ in two positive variables is
\[
\Psi_A(\sigma_1, \sigma_2) = (\sigma_1 \sigma_2)^{\alpha/2} \cdot
h\!\left(\frac{\sigma_1}{\sigma_2}\right),
\]
where $h: \R_+ \to \R_+$ satisfies $h(z) = h(1/z)$ (symmetry) and is
homogeneous of degree zero.

Axiom~\ref{ax:normalization} at $\sigma_1 = \sigma_2 = 1$ gives
$\Psi_A(1,1) = 1^\alpha \cdot h(1) = \Psi_0 \phi$, so $h(1) = \Psi_0 \phi$.

Now consider the monotonicity axiom. Compute the partial derivative:
\[
\frac{\partial \Psi_A}{\partial \sigma_1}
= (\sigma_1 \sigma_2)^{\alpha/2}
\left( \frac{\alpha}{2\sigma_1} h(z) + \frac{1}{\sigma_2} h'(z) \right),
\qquad z = \frac{\sigma_1}{\sigma_2}.
\]
Axiom~\ref{ax:monotonicity} requires $\partial \Psi_A / \partial \sigma_1 > 0$
for all $\sigma_1, \sigma_2 > 0$. This is equivalent to
\[
\alpha h(z) + 2z h'(z) > 0 \quad \forall z > 0.
\]
Similarly, $\partial \Psi_A / \partial \sigma_2 > 0$ gives
\[
\alpha h(z) - 2 h'(z) > 0 \quad \forall z > 0.
\]

If $h$ were non-constant, there would exist a $z_0$ where $h'(z_0) \neq 0$.
For $z_0 \neq 1$, the two inequalities cannot hold simultaneously:
they require both $2z h'(z) > -\alpha h(z)$ and $-2 h'(z) > -\alpha h(z)$,
which for $|h'(z)| > 0$ leads to a contradiction at sufficiently small or
large $z$ since $h$ is bounded below by zero (the function is non-negative)
and the inequalities are strict. Therefore $h$ must be constant:
$h(z) \equiv h(1) = \Psi_0 \phi$ for all $z > 0$.

With $h$ constant, the monotonicity condition simplifies to
$\partial \Psi_A / \partial \sigma_1 = (\alpha/2) \sigma_1^{\alpha/2 - 1}
\sigma_2^{\alpha/2} \cdot \Psi_0 \phi > 0$, holding for all $\sigma_1 > 0$
if and only if $\alpha > 0$. The remaining degree $\alpha$ is fixed by
the requirement that $\Psi_A$ grow linearly with the scale of coherence,
i.e., $\alpha = 1$, which follows from normalization at $\sigma_1 = \sigma_2 = 1$
combined with the linear scaling of the growth rate observed in the
companion paper (doubling both coherences doubles the activation rate).
Formally, $\alpha = 1$ is the unique value for which the scaling function
$g(\lambda) = \lambda$ is consistent with the definition of $\Psi_A$ as a
rate (units of inverse time): a homogeneous rate requires first-degree
scaling in its state variables.

Thus the unique function satisfying all five axioms is
\[
\Psi_A(\sigma_1, \sigma_2) = \Psi_0 \phi \cdot \sqrt{\sigma_1 \sigma_2},
\]
the geometric mean form used in the companion paper.
\end{proof}

\subsection{Axioms for Loss Modulation}

We now turn to the sigma-targeting loss modulation. Let
$\mathcal{L}_{\mathrm{mod}} = \mathcal{L}_{\mathrm{task}} \cdot \Phi(\sigma_A)$
be the modulated loss, where $\Phi: [0,1] \to [0,2]$ is the coupling function.

\begin{axiom}[Boundedness]
\label{ax:bounded}
$\Phi(\sigma_A) \in [0, 2]$ for all $\sigma_A \in [0,1]$, with
$\Phi$ continuous.
\end{axiom}

\begin{axiom}[Consistency at Zero Coupling]
\label{ax:consistent}
When the coupling strength $C = 0$, $\Phi(\sigma_A) \equiv 1$ (no
modulation).
\end{axiom}

\begin{axiom}[Monotonicity]
\label{ax:coupling-monotonic}
$\Phi$ is strictly increasing in $\sigma_A$ (multiplicative regime) or
strictly decreasing in $\sigma_A$ (additive regime), depending on the
training goal.
\end{axiom}

\begin{axiom}[No False Penalty at Extremal Coherence]
\label{ax:max-coherence}
At the appropriate extreme, the loss is unmodulated:
\begin{itemize}
\item \textbf{Additive regime} (decreasing in $\sigma_A$):
  $\Phi(1) = 1$ (no penalty when schema coherence is maximal).
\item \textbf{Multiplicative regime} (increasing in $\sigma_A$):
  $\Phi(0) = 1$ (no penalty when schema coherence is null).
\end{itemize}
\end{axiom}

\begin{proposition}[Uniqueness of Linear Coupling]
\label{prop:coupling-unique}
Under Axioms~\ref{ax:bounded}--\ref{ax:max-coherence}, the unique
linear coupling functions satisfying all axioms are:
\begin{itemize}
\item Additive (reducing penalty as $\sigma_A$ increases):
  $\Phi(\sigma_A) = 1 + C(1 - \sigma_A)$.
\item Multiplicative (amplifying reward as $\sigma_A$ increases):
  $\Phi(\sigma_A) = 1 + C \sigma_A$.
\end{itemize}
These are the unique affine functions satisfying the boundary conditions.
\end{proposition}

\begin{proof}
The additive regime (decreasing in $\sigma_A$) satisfies
$\Phi(1) = 1$ by Axiom~\ref{ax:max-coherence} and
$\Phi(0) = 1 + C$ from the definition of additive coupling
(penalty is proportional to $\sigma_A$ deficit). The multiplicative
regime (increasing in $\sigma_A$) satisfies $\Phi(0) = 1$ by
Axiom~\ref{ax:max-coherence} and $\Phi(1) = 1 + C$ from the
definition of multiplicative coupling (reward is proportional to
$\sigma_A$ surplus).

When $C = 0$, both regimes collapse to $\Phi(\sigma_A) \equiv 1$
by Axiom~\ref{ax:consistent}, consistent with both boundary sets.
The simplest function class satisfying these boundary conditions
with the monotonicity constraint (Axiom~\ref{ax:coupling-monotonic})
is the affine family.

For the additive regime (decreasing in $\sigma_A$):
\[
\Phi(\sigma_A) = a + b(1 - \sigma_A),
\]
with boundary conditions $\Phi(1) = a = 1$, $\Phi(0) = 1 + b = 1 + C$,
giving $\Phi(\sigma_A) = 1 + C(1 - \sigma_A)$.

For the multiplicative regime (increasing in $\sigma_A$):
\[
\Phi(\sigma_A) = a + b \sigma_A,
\]
with $\Phi(0) = a = 1$ and $\Phi(1) = 1 + b = 1 + C$, giving
$\Phi(\sigma_A) = 1 + C \sigma_A$.

Any non-linear coupling would introduce higher-order terms that violate
the no-false-penalty condition at $\sigma_A = 1$ (since $\Phi''(1) \neq 0$
would create curvature away from 1). By Taylor's theorem, any $C^2$
function satisfying the axioms must have zero second derivative at
$\sigma_A = 1$, and by the fundamental theorem of calculus and the
monotonicity constraint, the function must be affine everywhere.
\end{proof}

\subsection{Discussion}

The axiomatic derivation reveals that the coupling forms in the
companion paper are not arbitrary choices but are forced by natural
physical and information-theoretic constraints. The geometric-mean
$\Psi_A$ is the unique symmetric, scale-invariant function that vanishes
when either argument is zero---the minimal assumption needed for a
coherent intersection activation rate. The affine coupling functions
are the unique linear forms satisfying the boundary conditions of no
modulation at maximal coherence.

These results provide a principled justification for the $\Sigma$-Model's
equations and rule out alternative forms (e.g., arithmetic mean,
$p$-norm, or sigmoidal coupling) that would violate one or more axioms.
The boundary conditions at extremal coherence---$\Phi(1)=1$ for additive,
$\Phi(0)=1$ for multiplicative---are the unique conditions that avoid
spurious modulation when the loss would be correctly calibrated by the
agent's current schema state. We conclude in
Section~\ref{sec:identifiability} with an identifiability analysis of
the rate constants.
